% =====================================================================
%  starter-article.tex — copy this to start a new ARTICLE or fatwā reply.
%
%  An article is a book with the cover and the table of contents switched
%  off: it opens straight into a running head on page 1, and leans on
%  footnotes.
%
%  BUILD — from the kit's main folder:   bash tools/build.sh starter-article
%
%  THE ONE RULE THAT IS NOT TYPOGRAPHIC
%    Nothing here inserts content you did not type.
% =====================================================================
\documentclass[12pt]{article}
\newcommand{\ibfontpath}{fonts/}
\usepackage[article]{islamicbooklet}

\begin{document}

\bookstyle{Short title for the running head}

\doctitle{The Full Title of the Article}

Opening prose. There is no cover: page 1 is body text under a running
head.

% ---------------------------------------------------------------------
%  A MARGIN LABEL, pointing at the block beside it, and two shaded
%  blocks — e.g. a question as posed, and the answer given.
% ---------------------------------------------------------------------
\sidenote{The question reads:}
\begin{shaded}{shadequestion}
The question as posed, on a pale straw ground.
\end{shaded}

\begin{shaded}{shadeanswer}
The answer, on a pale green ground.
\end{shaded}

% ---------------------------------------------------------------------
%  THE COLOUR KEY. Write it out by hand: it changes per document, and an
%  automatic version would list colours you never used.
% ---------------------------------------------------------------------
\unheadingorheading{How to read this document}

\begin{itemize}
  \item \inquran{Dark red} is for the āyāt of the Qurʾān and their
        meanings in English.
  \item \inhadith{Green} is for the aḥādīth.
  \item \inathar{Blue} is for a saying of other than the Prophet
        \saw{} that has been quoted.
  \item \g{Grey round brackets} clarify the meaning of a word;
        \ed{grey square brackets} are the editor's note.
\end{itemize}

\divider

Ordinary prose continues. Footnotes are continuous through the document
and links inside them are live.\footnote{Taken from
\link{https://example.com}{this source}.}

\end{document}
